\documentclass[12pt,a4paper]{article}
\newcommand{\CadernoDisciplina}{Análise Exploratória de Dados}
\newcommand{\CadernoCorHex}{17365D}
% Preâmbulo compartilhado dos cadernos de exercícios UNIFACOL.
% Definir \CadernoDisciplina e, opcionalmente, \CadernoCorHex antes deste input.
\providecommand{\CadernoDisciplina}{Disciplina}
\providecommand{\CadernoCorHex}{17365D}
\usepackage[a4paper,margin=2.05cm,headheight=15pt]{geometry}
\usepackage[T1]{fontenc}
\usepackage[utf8]{inputenc}
\usepackage[brazil]{babel}
\usepackage{lmodern,microtype,amsmath,amssymb}
\usepackage{enumitem,booktabs,tabularx,array,longtable,adjustbox,multirow}
\usepackage{xcolor,hyperref,graphicx,fancyhdr}
\usepackage[most]{tcolorbox}
\usepackage{tikz}
\usetikzlibrary{arrows.meta,positioning,fit,calc,patterns,shapes.geometric}
\definecolor{disciplinecolor}{HTML}{\CadernoCorHex}
\definecolor{stimulusgrey}{HTML}{EEF2F5}
\definecolor{answerorange}{HTML}{FFF0DE}
\definecolor{answerframe}{HTML}{B85C00}
\definecolor{supportgrey}{HTML}{F7F7F7}
\hypersetup{colorlinks=true,linkcolor=disciplinecolor,urlcolor=disciplinecolor}
\newcolumntype{Y}{>{\raggedright\arraybackslash}X}
\newcolumntype{C}{>{\centering\arraybackslash}X}
\setlength{\parindent}{0pt}
\setlength{\parskip}{.42em}
\setlist[enumerate]{leftmargin=*,itemsep=.28em,topsep=.35em}
\pagestyle{fancy}
\fancyhf{}
\lhead{UNIFACOL --- \CadernoDisciplina}
\rhead{Caderno ENADE --- 2026.2}
\cfoot{\thepage}
\newcommand{\ItemHeader}[4]{%
  \par\medskip\noindent
  \colorbox{disciplinecolor}{\parbox{\dimexpr\textwidth-2\fboxsep\relax}{\color{white}\bfseries Questão #1\hfill #2}}%
  \par\smallskip\noindent\textbf{Competência:} #3\par
  \noindent\textbf{Suporte:} #4\par\smallskip
}
\newcommand{\TextoBase}{\noindent\colorbox{stimulusgrey}{\parbox{\dimexpr\textwidth-2\fboxsep\relax}{\textbf{TEXTO-BASE}}}\par\smallskip}
\newcommand{\Comando}{\par\smallskip\noindent\textbf{COMANDO.} }
\newcommand{\FonteDidatica}{\par\smallskip\noindent{\footnotesize\textit{Fonte: situação e dados elaborados para fins didáticos.}}}
\newcommand{\FonteExterna}[1]{\par\smallskip\noindent{\footnotesize\textit{Fonte: #1}}}
\newcommand{\EspacoDiscursiva}{\par\noindent\rule{\textwidth}{.2pt}\vspace{1.15cm}\par\noindent\rule{\textwidth}{.2pt}}
\newtcolorbox{AnswerBox}{enhanced,breakable,colback=answerorange,colframe=answerframe,boxrule=.7pt,arc=1mm,title=Resposta explicada,fonttitle=\bfseries,coltitle=white,before skip=7pt,after skip=7pt}
\newtcolorbox{DocumentBox}[1][Documento ou artefato]{enhanced,breakable,colback=supportgrey,colframe=black!55,boxrule=.6pt,arc=.7mm,title={#1},fonttitle=\bfseries,coltitle=white,before skip=7pt,after skip=7pt}
\newtcolorbox{MemoBox}{enhanced,breakable,colback=white,colframe=disciplinecolor,boxrule=.7pt,arc=.7mm,title=Memorando,fonttitle=\bfseries,coltitle=white,before skip=7pt,after skip=7pt}
\newcommand{\LegendaDidatica}[1]{\par\smallskip{\footnotesize\textit{#1. Fonte: elaborado para fins didáticos.}}}

\setcounter{secnumdepth}{0}
\begin{document}
\begin{center}
{\Large\bfseries UNIFACOL}\\[2mm]
{\large Segunda Chamada}\\[1mm]
{\Large\bfseries VERSÃO B}
\end{center}
\begin{DocumentBox}[Identificação]
\begin{tabularx}{\textwidth}{@{}p{.48\textwidth}p{.48\textwidth}@{}}
Estudante: \hrulefill & Turma: \hrulefill\\[3mm]
Professor: Cayo Oliveira & Data: \rule{2.7cm}{.2pt}/2026
\end{tabularx}
\end{DocumentBox}
\textbf{Instruções.} Esta avaliação contém seis questões objetivas e duas discursivas. Nas objetivas, assinale uma única alternativa. Nas discursivas, mostre cálculos intermediários e justifique decisões com as evidências do texto-base. Respostas sem desenvolvimento podem não receber pontuação integral.
\bigskip
% origem: AED-C11-O01
\ItemHeader{1}{Objetiva}{Selecionar ferramenta pela fronteira de governança}{Consulta executiva multifuente}
\TextoBase
Uma rede varejista possui métricas certificadas no warehouse, políticas por linha e um catálogo corporativo. O protótipo com respostas em linguagem natural será usado por 300 gestores. A diretoria exige que toda resposta identifique métrica, filtros, consulta executada e fonte; protótipos fora da camada semântica não podem publicar números oficiais.
A escolha será submetida ao comitê de dados, que rejeita soluções incapazes de reconstruir uma resposta contestada.
\FonteDidatica
\Comando A decisão arquitetural mais coerente é
\begin{enumerate}[label=\Alph*)]
  \item escolher a interface com o gráfico mais elaborado e recriar métricas no prompt.
  \item permitir SQL livre para todos os gestores e auditar apenas respostas denunciadas como erradas.
  \item usar qualquer chatbot, pois o catálogo torna desnecessário aplicar políticas na consulta.
  \item exportar dados para planilhas individuais e comparar as respostas manualmente ao final do mês.
  \item integrar a experiência conversacional à camada semântica e à identidade corporativa, registrando consulta, versão e evidência.
\end{enumerate}

% origem: AED-C10-O09
\ItemHeader{2}{Objetiva}{recusa}{Situação profissional e suporte quantitativo}

\TextoBase

A pergunta ``qual cliente fraudará amanhã?'' não possui rótulo, modelo validado nem base individual autorizada. Produzir uma lista nominal criaria inferência de alto impacto sem evidência. O agente deve recusar de forma útil e oferecer análise agregada permitida.

\FonteDidatica

\Comando Qual resposta do agente é compatível com as evidências e com o risco da solicitação?

\begin{enumerate}[label=\Alph*)]
  \item criar rótulo no prompt
  \item recusar inferência e explicar dados/competência ausentes, sem fabricar SQL
  \item retornar todos os clientes
  \item gerar ranking por gasto
  \item usar receita como fraude
\end{enumerate}

% origem: AED-C07-O09
\ItemHeader{3}{Objetiva}{recomendação}{Situação profissional e suporte quantitativo}

\TextoBase

Um campus adotou intervenção de permanência no mesmo mês em que o calendário acadêmico mudou; os demais não adotaram. A evasão caiu no campus participante. A secretaria precisa decidir o próximo passo sem atribuir toda a mudança à intervenção por simples comparação antes/depois.

\FonteDidatica

\Comando Selecione o desenho de continuidade que gera evidência comparável antes de uma expansão definitiva.

\begin{enumerate}[label=\Alph*)]
  \item comparar apenas antes/depois
  \item ocultar mudança de calendário
  \item recomendar piloto ampliado com comparação e métricas prévias, não expansão causal definitiva
  \item cancelar porque não há ensaio perfeito
  \item expandir imediatamente por correlação temporal
\end{enumerate}

% origem: AED-C07-D11
\ItemHeader{4}{Discursiva}{Integrar evidências e justificar decisão}{Caso profissional do capítulo}

\TextoBase

Uma base de evasão deveria possuir uma linha por estudante-semestre, mas transferências duplicam matrículas, frequência apresenta falhas por campus e regras mudaram durante o período. A recomendação final será usada para priorizar atendimento e precisa sobreviver a auditoria.

\FonteDidatica

\Comando Desenhe, em até 15 linhas, o fluxo da fonte à recomendação. Inclua teste de grão, ausência e domínio, transformação versionada, visual apropriado, reconciliação e evidência necessária para auditoria.
\EspacoDiscursiva

% origem: AED-C02-O07
\ItemHeader{5}{Objetiva}{Distinguir variância populacional e amostral}{Equipamentos}
\TextoBase Um laboratório registrou as idades $2,4,4,6$ anos de quatro equipamentos. A média é 4 e a soma dos desvios quadráticos é 8. Em um relatório, os quatro equipamentos constituem toda a sala; em outro, são tratados como amostra de equipamentos da rede. A escolha do divisor deve acompanhar essas duas finalidades.
\FonteDidatica
\Comando Se forem população e, alternativamente, amostra, as variâncias são
\begin{enumerate}[label=\Alph*)]
  \item 0 e 8, porque os desvios somam zero.
  \item $\sqrt8$ e $\sqrt3$.
  \item $8/4=2$ e $8/3\approx2{,}67$.
  \item $8/2=4$ em ambos os casos.
  \item $8/3$ e $8/4$, porque amostras usam divisor maior.
\end{enumerate}

% origem: AED-C13-O09
\ItemHeader{6}{Objetiva}{Aplicar recusa segura}{Pergunta sem evidência autorizada}
\TextoBase
O usuário pede previsão de demissões por pessoa. O agente tem apenas dados agregados autorizados, e a política proíbe inferência individual. Uma ferramenta de busca não encontra fonte adicional.
A resposta será usada em decisão de crédito; ausência de fonte autorizada exige abstenção útil, sem fabricar aproximação nem executar ferramenta privilegiada.
\FonteDidatica
\Comando O comportamento correto é
\begin{enumerate}[label=\Alph*)]
  \item executar várias buscas até alguma resposta parecer confiante.
  \item guardar a pergunta na memória e responder depois sem nova autorização.
  \item recusar a inferência individual, explicar a restrição e oferecer análise agregada autorizada ou encaminhamento humano.
  \item estimar indivíduos a partir da média da unidade e omitir a incerteza.
  \item pedir ao modelo que invente dados sintéticos com nomes reais.
\end{enumerate}

% origem: AED-C11-D11
\ItemHeader{7}{Discursiva}{Projetar matriz de comparação governada}{Seleção institucional}
\TextoBase
Uma universidade quer escolher entre as seis abordagens do capítulo para perguntas acadêmicas e financeiras. Há métricas certificadas, dados pessoais, grupos com menos de cinco estudantes e orçamento limitado.
O piloto precisa incluir perguntas rotineiras, ambíguas e proibidas, além de perfis com permissões distintas e decisões de publicação observáveis.
\FonteDidatica
\Comando Proponha, em até 15 linhas, um piloto comparável. Defina perguntas, perfis, fontes, critérios de semântica, acesso, evidência, auditoria, custo e latência; explique como tratar grupos pequenos e formule uma regra de aprovação sem escolher marca antecipadamente.
\EspacoDiscursiva

% origem: AED-C03-O03
\ItemHeader{8}{Objetiva}{Calcular união sem dupla contagem}{Eventos em um dado}
\TextoBase Na validação de uma regra de eventos, utiliza-se um dado honesto com seis resultados equiprováveis. Define-se $A=\{2,4,6\}$, resultado par, e $B=\{5,6\}$, resultado maior que quatro. O resultado 6 pertence aos dois eventos, portanto a probabilidade da união precisa evitar dupla contagem.
\FonteDidatica
\Comando $P(A\cup B)$ é
\begin{enumerate}[label=\Alph*)]
  \item $2/3-1/2=1/6$.
  \item $(3+2-1)/6=4/6$.
  \item $3/6+2/6=5/6$.
  \item $6/6$, pois união inclui o espaço inteiro.
  \item $1/6$, pois somente 6 pertence aos dois.
\end{enumerate}

\end{document}
